% Einheit 4 von 4 – Zuordnungstypen erkennen und anwenden (bank/zuordnungen/e4.jsonl, zusammenbau v0.1)
%% TODO Zweigzeile Teil 1: was hier gelernt wird (Satz in Schülersprache aus dem Einheitstitel) – Einheit 4
\einheitenkopf[e4]{Einheit 4 von 4 · Zuordnungstypen erkennen und anwenden}
\zweigzeile{ab Kl. 7 · P10 oft}
\setcounter{aufgabe}{22}

%% TODO Titel: Ich-kann-Satz fehlt in der Bank (Platzhalter: Kettenname) – Nr. 23
%% TODO Anweisung: ein Satz über der ersten Teilaufgabe fehlt in der Bank – Nr. 23
\begin{aufgabe}{Nullwert prüfen}
\begin{teile}
\teil Taxi: 4 € Grundgebühr und 2 € je km. Kostet eine Fahrt von 0 km auch 0 €? \\ \kreuz{ja} \\ \kreuz{nein}
\end{teile}
\end{aufgabe}

%% TODO \verfahren{…}: Verfahrensüberschrift in Sachsprache fehlt in der Bank (2.3 b)
%% TODO Titel: Ich-kann-Satz fehlt in der Bank (Platzhalter: Kettenname) – Nr. 24
%% TODO Anweisung: ein Satz über der ersten Teilaufgabe fehlt in der Bank – Nr. 24
\begin{aufgabe}{Erkennen}
\begin{teile}
\teil Je mehr Hefte du kaufst, desto … bezahlst du. \\ \kreuz{mehr} \\ \kreuz{weniger} \\ \kreuz{nicht vorhersagbar}
\end{teile}
\begin{teile}
\teil Ist die Zuordnung in der Tabelle proportional? Prüfe die Quotienten y : x. \\ \kreuz{proportional} \\ \kreuz{nicht proportional}
\end{teile}
\begin{teile}
\teil Ist die Zuordnung in der Tabelle proportional? Prüfe die Quotienten y : x. \\ \kreuz{proportional} \\ \kreuz{nicht proportional}
\end{teile}
\begin{teile}
\teil Ist die Zuordnung in der Tabelle proportional? Prüfe die Quotienten y : x. \\ \kreuz{proportional} \\ \kreuz{nicht proportional}
\end{teile}
\begin{teile}
\teil Ist die Zuordnung in der Tabelle proportional? Prüfe die Quotienten y : x. \\ \kreuz{proportional} \\ \kreuz{nicht proportional}
\end{teile}
\begin{teile}
\teil Ist die Zuordnung in der Tabelle proportional? Prüfe die Quotienten y : x. \\ \kreuz{proportional} \\ \kreuz{nicht proportional}
\end{teile}
\swfrage{Was bleibt gleich, was ändert sich?}
\begin{teile}
\teil Ist die Zuordnung in der Tabelle antiproportional? Prüfe die Produkte x · y. \\ \kreuz{antiproportional} \\ \kreuz{nicht antiproportional}
\end{teile}
\begin{teile}
\teil Fahrradverleih: 3 € Grundgebühr und 2 € je Stunde. Ist Mietdauer → Preis proportional? \\ \kreuz{proportional} \\ \kreuz{nicht proportional}
\end{teile}
\begin{teile}
\teil Welcher Zuordnungstyp gehört zum Graphen? \\ \kreuz{proportional} \\ \kreuz{antiproportional} \\ \kreuz{keins von beiden}
\end{teile}
\begin{teile}
\teil Welcher Zuordnungstyp liegt in der Tabelle vor? Prüfe erst die Quotienten, dann die Produkte. \\ \kreuz{proportional} \\ \kreuz{antiproportional} \\ \kreuz{keins von beiden}
\end{teile}
\swz[3]{Bei einem Paketdienst kosten 2 kg 8 € und 4 kg 14 €. Ist Gewicht → Preis proportional? Entscheide und begründe.}{}
\end{aufgabe}

%% TODO \verfahren{…}: Verfahrensüberschrift in Sachsprache fehlt in der Bank (2.3 b)
%% TODO Titel: Ich-kann-Satz fehlt in der Bank (Platzhalter: Kettenname) – Nr. 25
%% TODO Anweisung: ein Satz über der ersten Teilaufgabe fehlt in der Bank – Nr. 25
\begin{aufgabe}{Rate}
\begin{teile}
\teil Auf dem Tankbeleg steht: 40 l, 1,79 € je Liter, gesamt 71,60 €. Welche Angabe ist die Rate?
\end{teile}
%% TODO Darstellung für schwach fehlt in der Bank (zuordnungen-e4-k3-s1-v1; Abschnitt „Für schwache Schüler“)
\swz{Äpfel kosten 2,50 € je kg. Was kosten 4 kg?}{}
%% TODO Darstellung für schwach fehlt in der Bank (zuordnungen-e4-k3-s1-v2; Abschnitt „Für schwache Schüler“)
\swz{Kabel kostet 3 € je m. Was kosten 7 m?}{}
%% TODO Darstellung für schwach fehlt in der Bank (zuordnungen-e4-k3-s1-v3; Abschnitt „Für schwache Schüler“)
\swz{Nachhilfe kostet 15 € je Stunde. Was kosten 6 Stunden?}{}
%% TODO Darstellung für schwach fehlt in der Bank (zuordnungen-e4-k3-s1-v4; Abschnitt „Für schwache Schüler“)
\swz{Parken kostet 1,50 € je Stunde. Was kosten 3 Stunden?}{}
%% TODO Darstellung für schwach fehlt in der Bank (zuordnungen-e4-k3-s1-v5; Abschnitt „Für schwache Schüler“)
\swz{Heizöl kostet 0,90 € je Liter. Was kosten 500 l?}{}
\swfrage{Was bleibt gleich, was ändert sich?}
%% TODO Darstellung für schwach fehlt in der Bank (zuordnungen-e4-k3-s2-v1; Abschnitt „Für schwache Schüler“)
\swz{Kartoffeln kosten 1,50 € je kg. Wie viele Kilogramm bekommt man für 6 €?}{}
%% TODO Darstellung für schwach fehlt in der Bank (zuordnungen-e4-k3-s3-v1; Abschnitt „Für schwache Schüler“)
\swz{Ein Radfahrer fährt 45 km in 3 Stunden. Wie schnell fährt er im Schnitt?}{}
%% TODO Darstellung für schwach fehlt in der Bank (zuordnungen-e4-k3-s4-v1; Abschnitt „Für schwache Schüler“)
\swz{Ein Bus fährt 150 km mit durchschnittlich 50 km/h. Wie lange dauert die Fahrt?}{}
%% TODO Darstellung für schwach fehlt in der Bank (zuordnungen-e4-k3-s5-v1; Abschnitt „Für schwache Schüler“)
\swz[3]{Ein Lieferwagen verbraucht 8 l Diesel auf 100 km und fährt im Jahr 30\,000 km. Ein Liter Diesel kostet 162,5 ct. Berechne die Dieselkosten für ein Jahr in Euro. \hfill (P10 2014 OS)}{}
\swz[3]{Eine Wandergruppe geht mit 4 km/h. Sie ist um 9:50 Uhr an einer Hütte, macht dort 30 Minuten Pause und geht dann 3\,000 m bis zum See. Wann kommt sie am See an? \hfill (P10 2014 OS)}{}
\end{aufgabe}

%% TODO Titel: Ich-kann-Satz fehlt in der Bank (Platzhalter: Kettenname) – Nr. 26
%% TODO Anweisung: ein Satz über der ersten Teilaufgabe fehlt in der Bank – Nr. 26
\begin{aufgabe}{Gegenbeispiele (Alter und Größe, Tarif mit Grundgebühr)}
%% TODO Darstellung für schwach fehlt in der Bank (zuordnungen-e4-k4-s1-v1; Abschnitt „Für schwache Schüler“)
\swz[3]{Nenne ein Beispiel für eine Zuordnung „je mehr …, desto mehr …“, die nicht proportional ist, und begründe kurz.}{}
\end{aufgabe}

%% TODO Titel: Ich-kann-Satz fehlt in der Bank (Platzhalter: Kettenname) – Nr. 27
%% TODO Anweisung: ein Satz über der ersten Teilaufgabe fehlt in der Bank – Nr. 27
\begin{aufgabe}{Erkennen – Fehler finden, Begründen, Darstellungswechsel, Anwendung}
\swz[3]{Ein Auto fährt 90 km in 1 Stunde. Lukas rechnet: „Für 180 km braucht es eine halbe Stunde, denn die Strecke ist doppelt so lang.“ Finde den Fehler und rechne richtig.}{}
\swfrage{Warum ist die Zuordnung Anzahl Personen → Preis pro Person für eine Pizza nicht proportional? Begründe.}
\swa{Ein Brunnen liefert 5 l Wasser je Minute. Zeichne den Graphen Zeit in min → Wasser in l für 0 bis 4 Minuten und nenne den Zuordnungstyp.}{\begin{ksys}[xmin=0,xmax=5,ymin=0,ymax=20,xlabel=Zeit in min,ylabel=Wasser in l]\end{ksys}}
\swz[3]{Familie Yilmaz fährt 280 km in den Urlaub. Sie fährt im Schnitt 80 km/h und will um 14:00 Uhr ankommen. Wann muss sie spätestens losfahren?}{}
\end{aufgabe}

\uebersichtskasten{Erkennen: Tabelle – alle Quotienten y : x gleich → proportional; alle Produkte x · y gleich → antiproportional; sonst keins von beiden. \\ Text: „je mehr, desto mehr“ reicht nicht – auch 0 → 0 muss passen (Grundgebühr: nein). \\ Graph: Gerade durch den Ursprung → proportional; fallende Kurve → antiproportional. \\ Rate: Geschwindigkeit = Weg : Zeit; Zeit = Weg : Geschwindigkeit; Kosten = Menge · Preis je Einheit. \\ 3 km in 12 min → 0,25 km je Minute \quad 96 m : 4 m/s = 24 s}
